\begin{Definition}\label{d1_1}A set $\Theta = (J_3, J_5, \alpha, \beta)$, where $\alpha>0$, $\beta\in\mathbb{R}$, $J_3$ is a Jacobi matrix and~$J_5$ is a semi-inf\/inite real symmetric f\/ive-diagonal matrix with positive numbers on the second subdiagonal, is said to be \textit{a Jacobi-type pencil $($of matrices$)$}.
\end{Definition}
From this def\/inition we see that matrices $J_3$ and $J_5$ have the following form
\begin{gather}
J_3 =
\left(
\begin{matrix}
b_0 & a_0 & 0 & 0 & 0 & \cdots\\
a_0 & b_1 & a_1 & 0 & 0 & \cdots\\
0 & a_1 & b_2 & a_2 & 0 &\cdots\\
\vdots & \vdots & \vdots & \ddots \end{matrix}
\right),\qquad a_k>0,\quad b_k\in\mathbb{R},\quad k\in\mathbb{Z}_+, \label{f1_5} \\
J_5 = \left(
\begin{matrix}
\alpha_0 & \beta_0 & \gamma_0 & 0 & 0 & 0 & \cdots\\
\beta_0 & \alpha_1 & \beta_1 & \gamma_1 & 0 & 0 & \cdots\\
\gamma_0 & \beta_1 & \alpha_2 & \beta_2 & \gamma_2 & 0 & \cdots\\
0 & \gamma_1 & \beta_2 & \alpha_3 & \beta_3 & \gamma_3 & \cdots \\
\vdots & \vdots & \vdots &\vdots & \ddots \end{matrix}
\right),\qquad \alpha_n,\beta_n\in\mathbb{R},\quad \gamma_n>0,\quad n\in\mathbb{Z}_+. \label{f1_10}
\end{gather}

With a Jacobi-type pencil of matrices $\Theta$ one associates a system of polynomials $\{ p_n(\lambda) \}_{n=0}^\infty$, such that
\begin{gather*} p_0(\lambda) = 1,\qquad p_1(\lambda) = \alpha\lambda + \beta, \end{gather*}
and
\begin{gather}\label{f1_20}
(J_5 - \lambda J_3) \vec p(\lambda) = 0,
\end{gather}
where $\vec p(\lambda) = (p_0(\lambda), p_1(\lambda), p_2(\lambda),\dots)^{\rm T}$. Here the superscript ${\rm T}$ means the transposition. Polynomials $\{ p_n(\lambda) \}_{n=0}^\infty$ are said to be \textit{associated to the Jacobi-type pencil of matrices~$\Theta$}.

Observe that for each system of orthonormal polynomials on the real line with $p_0=1$ one may choose $J_3$ to be the corresponding Jacobi matrix (which elements are the recurrence coef\/f\/icients), $J_5 = J_3^2$, and $\alpha$, $\beta$ being the coef\/f\/icients of~$p_1$ ($p_1(\lambda) = \alpha\lambda + \beta$).

One can rewrite relation~(\ref{f1_20}) in the scalar form
\begin{gather} \gamma_{n-2} p_{n-2}(\lambda) + (\beta_{n-1}-\lambda a_{n-1}) p_{n-1}(\lambda) + (\alpha_n-\lambda b_n) p_n(\lambda) \notag \\
\qquad{} + (\beta_n-\lambda a_n) p_{n+1}(\lambda) + \gamma_n p_{n+2}(\lambda) = 0,\qquad n\in\mathbb{Z}_+, \label{f1_30}
\end{gather}
where $p_{-2}(\lambda) = p_{-1}(\lambda) = 0$, $\gamma_{-2} = \gamma_{-1} = a_{-1} = \beta_{-1} = 0$.
